% Generated draft: requires mathematical and source review.
\documentclass[11pt]{article}
\usepackage[margin=25mm]{geometry}
\usepackage{fontspec}
\setmainfont{Latin Modern Roman}
\usepackage{amsmath,amssymb,mathtools,unicode-math}
\setmathfont{Latin Modern Math}
\usepackage{enumitem,needspace,xurl,hyperref}
\setlength{\parindent}{0pt}
\setlength{\parskip}{5pt}
\setlength{\emergencystretch}{3em}
\newcommand{\N}{\mathbb{N}}
\newcommand{\Z}{\mathbb{Z}}
\newcommand{\Q}{\mathbb{Q}}
\newcommand{\R}{\mathbb{R}}
\newcommand{\C}{\mathbb{C}}
\newcommand{\F}{\mathbb{F}}
\newcommand{\A}{\mathbb{A}}
\newcommand{\PP}{\mathbb{P}}
\newcommand{\E}{\mathbb{E}}
\newcommand{\eps}{\varepsilon}
\newcommand{\dd}{\,\mathrm{d}}
\begin{document}
\section*{UnsolvedMath ID: 30002866}
\textbf{Source problem code:} catalogue entry
\subsection{Definitions and mathematical statement}
Let \(d\geq 1\) and let \(A=\C[x_1,\ldots,x_d]\). An ideal projector is a \(\C\)-linear map \(P:A\to A\) satisfying \(P^2=P\) whose kernel is an ideal of \(A\). The finite-dimensional setting mentioned in the supplied background is interpreted as \(\dim_{\C}\operatorname{ran}P<\infty\).

For clarity, adopt the following conventional interpretation of the terminology that the record leaves undefined. Put \(V=\operatorname{ran}P\) and \(n=\dim_{\C}V\). A set \(X=\{\xi_1,\ldots,\xi_n\}\subset\C^d\) of distinct points is unisolvent for \(V\) if the evaluation map \(V\to\C^n\), \(v\mapsto(v(\xi_1),\ldots,v(\xi_n))\), is an isomorphism. Its Lagrange projector \(L_X:A\to V\) is defined by
\[
 (L_Xp)(\xi_j)=p(\xi_j)\qquad(p\in A,\ 1\leq j\leq n).
\]
In this interpretation, \(P\) is Hermite if there exist sets \(X_k\) unisolvent for \(V\) such that
\[
 \lim_{k\to\infty}L_{X_k}p=Pp\qquad\text{for every }p\in A,
\]
where convergence takes place in the finite-dimensional space \(V\). For \(n=0\), the empty interpolation set gives the zero projector.

For polynomials \(f_1,\ldots,f_d\in A\), write \(f_i^{\mathrm{top}}\) for the highest-degree homogeneous part of a nonzero \(f_i\), and set \(f_i^{\mathrm{top}}=0\) if \(f_i=0\). One usual meaning of the assertion that \(f=(f_1,\ldots,f_d):\C^d\to\C^d\) has no roots at infinity is
\[
 \bigl\{z\in\C^d\setminus\{0\}:f_i^{\mathrm{top}}(z)=0\text{ for every }1\leq i\leq d\bigr\}=\varnothing.
\]
This convention describes the common zeros at infinity of the separately homogenized defining equations.

\textbf{Supplied theorem assertion, under these interpretations.} For every finite-rank ideal projector \(P:A\to A\) and every \(f_1,\ldots,f_d\in A\),
\[
 \ker P=\langle f_1,\ldots,f_d\rangle
 \quad\text{and}\quad
 f\text{ has no roots at infinity}
 \quad\Longrightarrow\quad
 P\text{ is Hermite}.
\]
Here \(\langle f_1,\ldots,f_d\rangle=\{\sum_{i=1}^d a_i f_i:a_i\in A\}\).

\textbf{Unresolved qualifications.} The supplied sentence does not explicitly specify finite rank, the definition or topology of Hermite approximation, or the precise meaning of roots at infinity. The definitions above make a conventional interpretation explicit; their agreement with the cited passage requires mathematical review. In particular, the homogenized equations need not describe the projective closure of their affine zero set without further qualification. The record's verification field identifies the sentence as a proved sufficient condition, rather than an open problem, and supplies no separate recoverable conjecture. The broader Hermite-projector conjecture mentioned in the background is therefore not substituted for this assertion.
\subsection{Short English statement}
The supplied assertion says that an ideal projector whose kernel has \(d\) polynomial generators defining a system with no roots at infinity is Hermite. A finite-rank interpretation and explicit conventions for the undefined terms are given above, subject to review. The record identifies this passage as a theorem extraction and does not provide a separate precise open problem.
\subsection{Sources}
\begin{itemize}
\item Local UnsolvedMath registry: supplied complete record for permanent ID \(30002866\), source code \texttt{catalogue entry}, titled \emph{Hermite-Projector Theorem Extraction}. No proposer is supplied; the proposed year is recorded as 2015.
\item Supplied report citation: \emph{Multivariate Splines and Algebraic Geometry (2016)}, Oberwolfach Reports, pp.\ 1139--1200, DOI \texttt{10.4171/owr/2015/21}. \url{https://doi.org/10.4171/owr/2015/21}. The literature-source list labels this reference \emph{Oberwolfach Report 21}. The supplied background describes Oberwolfach Reports as documenting mathematical workshops and their open-problem discussions.
\item Supplied literature reference: Shekhtman, \emph{On One Class of Hermite Projectors}. \url{https://doi.org/10.1007/s00365-015-9311-5}. No more detailed theorem attribution or location is supplied.
\item Literature assessment dated exactly \texttt{2026-08-21}: supplied dated curation, not independently verified. It labels the record \texttt{partially\_solved} and describes the supplied theorem as a proved partial result supporting the Hermite-projector conjecture that does not cover all finite-dimensional ideal projectors. The background explicitly calls this dated literature triage derived from the supplied OWR curation and requests independent verification of time-sensitive status and cited claims.
\item The supplied statement status is \texttt{unrecoverable}. Its verification note identifies the passage as a proved sufficient condition, not an open problem, states that no separate clean statement is promoted, and retains the original record for provenance. These are supplied editorial assessments, not independent mathematical verification.
\end{itemize}
\end{document}
